\section{Discussion}
\label{sec:discussion}

\subsection{What the sampling comparison establishes}

The comparison against the matched-budget baselines turns on the record budget, and on which error
statistic is read. Under aggressive compression, where the budget is too small for a random draw to span
the target system, informed $L_0$ selection holds the lowest error because the gates spend the budget on
the records the targets constrain. As the budget grows the methods converge: on the headline median the
edge passes to random + reweight once almost any subset can be reweighted to fit (near 5{,}000 records
here), while on the tail-sensitive mean informed $L_0$ stays ahead at every budget, because it removes the
seed-to-seed instability that swings a small random or survey draw's mean by thousands of per cent
(Section~\ref{sec:results}). The result is graceful degradation under compression rather than a blanket
accuracy win, and the budget at which the median crossover happens is itself a finding. Past it the case
for informed $L_0$ rests on its controls, on its near-zero generalization gap, and on the tail it keeps
bounded, rather than on a median accuracy gap. It removes the risk of a poor random draw and sets the
dataset size, the geographic emphasis, and the weight concentration through explicit parameters in one
optimization. That control matters most in the build-big-then-prune setting the method is designed for,
where a pipeline assembles a candidate universe far larger than it can ship and must reduce it to a fixed
budget under subnational calibration targets; this study tests the national and state surface, and the
subnational case is left to future work. Reporting out-of-sample error on held-out families keeps the
comparison honest, since any method matches the targets it was fit to more tightly than ones it never saw.

The convex $L_1$ arm sharpens the contrast. $L_1$ pursues the same target-informed sparse selection, but
through a single penalty that must both decide which records to drop and how much to shrink the survivors. Forcing roughly two thousand of seventy-five thousand records to survive demands a
large penalty, which soft-thresholds the retained weights toward zero and collapses the calibration fit:
its out-of-sample median ARE is pinned near 100\% at the tightest budgets and its mean runs to several tens
of thousands of per cent, recovering toward informed $L_0$ only once the budget is loose enough that little
shrinkage is needed (Figure~\ref{fig:budget_frontier}). Informed $L_0$ avoids this because the Hard
Concrete gate and the fitted weight are separate variables: the gate decides membership and the weight is
free to take whatever positive value the targets require. Decoupling selection from shrinkage is what lets
$L_0$ degrade gracefully where the convex relaxation does not, and it is the core of the $L_0$-versus-$L_1$
comparison. $L_1$ is more stable to optimize and can be the better choice where sparsity is not forced
hard; the case for $L_0$ here rests on the joint selection-and-weighting framing and the explicit control
over the retained count, not on a claim that it is the most accurate sparse method in every setting.

\subsection{What the controls establish}

The second contribution is operational: the same framework produces datasets of different sizes by
changing the $L_0$ penalty, focuses on a geographic level by weighting its targets, and shapes the
weight distribution through two concentration controls, a soft $L_2$ penalty and a hard per-record cap.
We treat the last of these as a measured frontier rather than a fixed setting: sweeping the $L_2$ penalty
traces what tightening concentration (raising the effective sample size and lowering the largest weights)
costs in calibration accuracy, so an operator can choose a point on that trade-off rather than accept
whatever concentration the unconstrained fit produces. A random-sampling baseline
can also produce a dataset of a chosen size, but it cannot make the choice of records respond to the
targets, and it controls concentration only through the reweighting step. The value here is that size,
geographic focus, and weight stability are set through explicit parameters in one optimization rather
than through separate heuristics.

\subsection{Trade-offs and limitations}

The method does not remove modelling judgment, and several choices stay consequential; the optimizer's
hyperparameters are one such choice, reported in full in Appendix~\ref{app:hyperparameters}.

\subsubsection{Target selection}

The target set is a prioritized subset, not a complete inventory. The reasons for including each
family are stated in Appendix~\ref{app:targets}, and the empirical claims are tied to the exact
targets in each run's manifest. A different target set could change the comparison, and the paper
does not claim the chosen set is maximal.

\subsubsection{When informed selection helps}

Informed selection should help most when the targets are specific and local, so the choice of
records matters. When the targets are few and broad, a random draw may already contain the records
needed to match them, and the advantage of informed selection narrows. The paper reports the regimes
it tests and does not extrapolate beyond them.

\subsubsection{Weight concentration and effective sample size}

Matching a demanding fiscal target system from a uniform prior can require the optimizer to load a
large share of the population onto a small number of records, so the effective sample size can fall
well below the retained-record count even when the calibration fit is good. That concentration is
acceptable for the aggregates the dataset is calibrated to, but it is a real cost for any downstream
estimand that is sensitive to the weight distribution. We expose the cost rather than hide it: the
effective sample size and the largest weights are reported as primary results, and the $L_2$ sweep
shows how much accuracy raising the effective sample size costs. That cost depends on the budget. At
40{,}000 records $\lambda_{L_2}=10^{-4}$ raises the effective sample size by about seventy per cent (from
3{,}803 to 6{,}482) for under two points of median accuracy, while at the tightest budget it inflates the
error tail instead. The penalty is worth applying at large budgets and not at tight ones, which is why we
report the un-penalised frontier as the accuracy result and the penalty as a separate operability setting.

\subsubsection{Generalization across target families}

Holding out whole families is a deliberately hard test of generalization, harder than holding out a
random subset of targets, and it need not flatter any method. On the fixed holdout, informed $L_0$
generalizes with almost no gap between its in-sample and out-of-sample median error, while random +
reweight fits the in-sample targets far more tightly but pays about a thirty-point gap on families it
never saw; informed selection trades in-sample fit for a sample that carries over. The stricter family
rotation is less forgiving. On the target-weighted mean across folds the ranking favours the baselines,
but the gap is confined to the fold that holds out the dominant \texttt{irs\_soi} family: with 71\% of the
targets removed, every method's mean is dominated by a few near-zero-denominator targets, and the fold
medians stay close. We read the rotation as a limitation to be reported per fold rather than as a single
headline number, and widening the held-out surface further is left to future work.

\subsubsection{Degenerate targets and the choice of error metric}

A handful of targets have denominators small enough that one record's placement swings their relative
error past 100 per cent, which inflates the mean absolute relative error without reflecting calibration
quality. Rather than winsorize or trim the distribution by an arbitrary cutoff, we lead with the median,
name the affected targets in a one-time audit, and report the mean recomputed with exactly those named
targets removed. This keeps the metric honest and the dropped targets visible and attributable, but it
does mean the headline accuracy figure depends on a stated, auditable choice rather than on the raw mean
alone.

\subsubsection{Computational cost}

Building the calibration matrix and running gradient-based optimization cost more than a small
classical calibration, and informed $L_0$ adds the eight-step budget bisection on top of a single fit;
Appendix~\ref{app:cost} reports the runtime each method reaches against the accuracy it buys. That cost
is justified only when the problem requires it, which is why the comparison against a simpler
random-sampling baseline matters.

\subsection{Generalizability}

The data engineering here is specific to the United States, but the question is general. Many
microsimulation projects combine sources into a candidate dataset larger than they can ship, and then
have to reduce it. Where the reduction has to respect calibration targets and keep weights usable,
framing it as informed sampling, rather than as a random draw followed by reweighting, may carry over
to other countries and other engines. \populace{} makes that portability concrete, since the method
is a configuration choice rather than a fixed pipeline.
